\section{Introducción}

Esta práctica busca reforzar en el estudiante los conocimientos relacionados con el comportamiento del circuito magnético de dos tipos de núcleos ferromagnéticos: macizo y laminado. Se plantea el estudio de un circuito magnético, así como la medición del correspondiente Ciclo de Histéresis.

\section{Objetivos}

\subsection*{Objetivo General}
\begin{itemize}
    \item Analizar y comparar el comportamiento de núcleos ferromagnéticos macizo y laminado, evaluando sus curvas B-H, ciclo de histéresis y pérdidas asociadas.
\end{itemize}

\subsection*{Objetivos Específicos}
\begin{itemize}
    \item Comprender el ciclo de histéresis y sus parámetros característicos ($B_{\text{sat}}$, $B_r$, $H_c$).
    \item Determinar la curva B-H de un material ferromagnético.
    \item Calcular pérdidas por histéresis y corrientes parásitas (Foucault).
    \item Comparar las formas de onda de corriente y las pérdidas entre un núcleo macizo y uno laminado.
    \item Evaluar el efecto del laminado en la reducción de pérdidas por corrientes parásitas.
\end{itemize}
